\section{Discussion}

The two benefits of our design interact with the memory hierarchy differently.
For offloading, the main question is how much of the endpoint derivation logic can be lightweighted and placed in the HBM base die. We assume that the base die contains only a limited subset of the required hash/interleaving logic, while detailed placement information is supplied through metadata hints. This enables metadata-driven offloading without fully replicating the GPU's placement logic.
In contrast, the bisection bandwidth benefit is more closely related to caching organization. Our evaluation assumes that L2 cache blocks can be placed flexibly across the two GPU  dies depending on the transfer path. This departs from the tightly coupled L2 slice and memory controller structure of conventional GPUs, and may require support mechanisms similar to prior local/remote caching~\cite{cache_multichiplet_1, cache_multichiplet_2} and selective replication design~\cite{cache_multichiplet_3, cache_multichiplet_4}.
More broadly, once the memory controller is no longer tightly co-located with the on-die L2 cache and is instead connected through D2D links, revisiting the cache organization itself becomes a natural architectural consideration.

%Our two contributions interact with the legacy GPU memory system in different ways. The bisection-bandwidth benefit is primarily depends on caching semantics, whereas DME-based movement offloading additionally depends on how source and destination endpoints are derived under fine-grained address hashing. 
%In our baseline, we preserve the conventional home-slice organization in which ordinary cacheable requests are mapped to a home L2 slice or memory partition by the GPU's existing address-to-partition function. Under this organization, the HBM-side fabric operates below the home partition point and therefore does not directly change how ordinary cacheable accesss select L2 slices or participate in coherence. 
%As a result, the immediate beneficiaries of path diversity are DME traffic and software-designated non-coherent (NC) or explicitly managed read-mostly regions, where the fabric can be exploited without changing the legacy L2 slice and memory-partition contract. Extending this benefit to ordinary cacheable traffic would naturally motivate fabric-aware caching mechanisms, such as selective read-only/read-mostly replication across L2 slices~\cite{SelRepCache} or additional memory-side tag/directory support.

%The offloading interface follows the same boundary. Bulk movement is invoked explicitly through a runtime API, and the moved regions are assumed to be NC, so correctness does not depend on transient cache state. A natural extension is to let the compiler or runtime identify copy-like streams and lower them into instruction-level move requests served by the DME. Supporting cacheable regions would then require coordination with cache state and ordering, so that the system can choose between offloading, explicit cache management, and the conventional read/wirte path.

%This conservative boundary also shapes the offloading interface. In the evaluated design, bulk movement is invoked explicitly through a runtime API and the moved regions are assumed to be NC, so correctness does not depend on transient cache state. A more aggressive extension would be to let the compiler/runtime detect copy-like load/store streams and lower them into instruction-level move requests served by the DME. Once cacheable regions are allowed, however, the decision to offload can no longer ignore cache state: if the source or destination may already be present, especially dirty, in on-die caches, blindly delegating the operation to the DME risks stale data or ordering violations. A practical design would therefore need a cache-state-aware policy that either offloads the request, performs explicit cache management, or falls back to the conventional read/write path.

%Hashing affects the two contributions asymmetrically. The bisection bandwidth benefit does not require the HBM base die to replicate GPU placement logic. Once a request is issued, the fabric can still exploit multiple internal paths. For offloading, however, fine-grained interleaving can increase GPU-side command generation overhead if the GPU must materialize fully resolved per-chunk source and destination endpoints. This motivates a metadata-driven interface in which the GPU provides compact movement metadata and placement hints, while the base die reconstructs executable descriptors locally by implementing only the subset of hashing/interleaving logic needed for endpoint derivation.
%Such an interface could reduce GPU-side control overhead and control traffic while increasing DME autonomy for repeated or structured movement patterns. We view fabric-aware caching and metadata-driven endpoint derivation as promising extensions beyond the evaluated baseline.

%Hashing affects the two contributions asymmetrically. The bisection bandwidth argument does not require the HBM base die to replicate GPU placement logic; once a request is issued, the fabric can still exploit multiple internal paths. Offloading, in contrast, becomes less attractive if the GPU must materialize fully resolved per-chunk source/destination descriptors under fine-grained interleaving used to mitigate memory camping. A natural extension is a metadata-driven interface in which the GPU sends compact movement metadata and placement hints, while the base die reconstructs executable descriptors locally by replicating only the limited subset of hashing/interleaving logic needed for endpoint derivation. This could reduce GPU-side control overhead and descriptor traffic while increasing DME autonomy for repeated or structured movement patterns, but it requires a compact and stable GPU--HBM contract for placement semantics. We view such fabric-aware caching and metadata-driven endpoint derivation as promising co-designed extensions beyond the conservative baseline studied here.
